\section{Introduction}\label{sec:introduction}

Disease early warning and surveillance (DEWS) systems underpin outbreak detection, reporting, and response across hierarchically organized health institutions---from community health posts and primary facilities through district, zonal, and regional offices to national public health authorities. In low- and middle-income countries (LMICs) such as Ethiopia, these systems confront persistent structural failures: data and process centralization that introduces single-point-of-failure risks, security vulnerabilities arising from the absence of tamper-evident record-keeping, traceability gaps across the reporting hierarchy, reliance on manual completeness verification, and interoperability barriers that fragment information flow across organizational boundaries.

Blockchain technology has been proposed as an architectural remedy, offering decentralization, immutability, cryptographic integrity, and transaction traceability.\cite{yaga2018blockchain} However, first-generation blockchain-health architectures---including prior work on DEWS---exhibit four categories of shortcoming that limit their scientific and practical contribution.

First, \emph{monolithic on-chain storage} requires every participating node to replicate the full surveillance payload, producing unsustainable storage growth, high gas costs, and bandwidth demands incompatible with LMIC infrastructure.\cite{kelly2019key} Second, \emph{narrow privacy models} equate immutability with confidentiality, ignoring metadata-level exposure on public or permissioned ledgers, overbroad access across health-system tiers, and the absence of selective disclosure or permission revocation.\cite{finlayson2019adversarial} Third, \emph{unargued consensus assumptions}---typically Proof of Authority adopted without comparative evaluation---leave validator governance, collusion resilience, and trust concentration unaddressed.\cite{bohme2015building} Fourth, \emph{resource-uniform participation models} assume that all facilities, including rural health posts with intermittent connectivity and underpowered devices, can operate as full blockchain peers, thereby excluding the very populations surveillance is designed to protect.

These limitations motivate a fundamental architectural shift: from full ledger replication to \emph{selective anchoring}, where encrypted surveillance payloads reside off-chain while the blockchain retains only cryptographic hashes, permissions, provenance records, alerts, and audit events.

This paper presents a hybrid, privacy-layered, hierarchical surveillance architecture designed for LMIC deployment contexts. The architecture reframes Ethereum from a bulk data repository to a coordination and proof layer, integrating the following novelty contributions:

\begin{enumerate}
\item A \emph{selective-anchoring architecture} with hierarchical event processing: community and facility nodes submit encrypted case data to off-chain storage (IPFS), district and regional layers batch and validate submissions, and smart contracts commit only hashes, access-control state, and audit artifacts on-chain---eliminating full ledger replication of surveillance payloads.
\item A \emph{layered privacy model} combining zero-knowledge verification for privacy-preserving validation, dynamic smart-contract-based access control with grant, revoke, and expire semantics, and encrypted off-chain records---replacing immutability-only privacy assumptions.
\item \emph{Edge and fog preprocessing} at peripheral tiers to reduce bandwidth, latency, and compute burden, enabling resource-constrained facilities to participate without operating as heavy validator nodes.
\item A \emph{comparative consensus framework} evaluating Proof of Authority, Delegated Proof of Stake, PBFT variants, and Raft-based alliance models against LMIC workload characteristics, trust topology, and infrastructure constraints---with explicit validator governance, appointment, and failure-handling specifications.
\item \emph{FHIR-aware interoperability workflows} supporting standards-compliant selective disclosure and cross-organizational data exchange within the privacy architecture.
\item A functional prototype validating the smart-contract coordination layer, evaluated across three dimensions: performance (block generation timing), usability (Software Usability Measurement Inventory with 26 domain professionals), and security (expanded threat model addressing validator collusion, metadata leakage, overbroad tier exposure, and equity-aware governance).
\end{enumerate}

The remainder of this paper is organized as follows. Section~\ref{sec:background} presents background context and related work. Section~\ref{sec:methodology} describes the research methodology. Section~\ref{sec:model} details the proposed architecture. Section~\ref{sec:implementation} presents the prototype implementation. Section~\ref{sec:evaluation} reports evaluation results. Section~\ref{sec:discussion} discusses implications and limitations. Section~\ref{sec:conclusion} concludes the paper and identifies future research directions.

%% ===== BACKGROUND AND RELATED WORK =====
